\documentclass[11pt]{article}

\usepackage[T1]{fontenc}
\usepackage[utf8]{inputenc}
\usepackage{lmodern}
\usepackage{amsmath,amssymb,amsthm,mathtools}
\usepackage[a4paper,margin=28mm]{geometry}
\usepackage{microtype}
\usepackage{xurl}
\usepackage[hidelinks]{hyperref}
\usepackage{enumitem}
\usepackage{booktabs}
\usepackage{array}

\newtheorem{theorem}{Theorem}[section]
\newtheorem{lemma}[theorem]{Lemma}
\newtheorem{proposition}[theorem]{Proposition}
\newtheorem{corollary}[theorem]{Corollary}
\newtheorem{definition}[theorem]{Definition}
\newtheorem{remark}[theorem]{Remark}

\newcommand{\N}{\mathbb N}
\newcommand{\WIN}{\ensuremath{\mathrm{WIN}}}
\newcommand{\LOSS}{\ensuremath{\mathrm{LOSS}}}
\newcommand{\DRAW}{\ensuremath{\mathrm{DRAW}}}
\newcommand{\R}{\mathcal R}
\newcommand{\PhiM}{\Phi}
\newcommand{\odd}{\operatorname{odd}}
\newcommand{\vTwo}{v_2}
\newcommand{\source}{\operatorname{src}}

\title{The Two-Player $3n\mathbin{\pm}1$ Game:\\
A Binary Normal Form and a Well-Founded No-Draw Proof}
\author{Grisha Pochuev\\
  \small\href{mailto:n_854@mail.ru}{n\_854@mail.ru}\\
  \small\url{https://github.com/Grisha-Pochuev/3n-plus-minus-1-game}}
\date{Corrected manuscript, 10 August 2026}

\begin{document}
\maketitle

\begin{center}
\small\textbf{Audit status.} This manuscript presents the repaired human proof after
an external audit of an earlier version. The two invalid shortcuts identified by
that audit remain withdrawn. Independent external re-audit of the repaired bridge
is still pending, and no end-to-end Lean formalization is claimed.
\end{center}

\begin{abstract}
Consider the impartial two-player game on positive odd integers in which a move
from $n>1$ chooses $3n-1$ or $3n+1$ and removes all powers of two; reaching $1$
wins immediately. Arbitrary play need not terminate, since $5\to7\to5$ is a
legal cycle. The problem is whether optimal play nevertheless has finite
remoteness from every start. We pass to a conjugated binary game with an
expanding move $A$ and a strictly contracting move $B$, introduce constant-tail
coordinates, an embedded source map, and finite proof-height tokens for certified
winning and losing positions. The global argument is a well-founded marked
normalization: source decreases and certified token replacements form the outer
rank, while a typed factor/obligation system is well-founded inside each fixed
outer fibre. An external audit found that an earlier entry argument incorrectly
identified a smaller canonical coefficient with a smaller coefficient source and
also used a reverse-factor lemma outside its valid exponent range. We repair both
issues. The key final step is a one-shot entry component that permits a returned
inner source to be larger than the original source exactly once, during typed
normalizer initialization, while forbidding any repeated unranked reset. A direct
phase analysis of the remaining factor-free exponent-two base entry closes the
last attachment case. The resulting human proof rules out all draw positions.
Machine regressions and symbolic certificates are supporting checks with a
strictly narrower scope than the human theorem.
\end{abstract}

\tableofcontents

\section{The game and the theorem}

The current state is a positive odd integer $n$. For $n>1$ the player to move
chooses one of
\[
  3n-1,\qquad 3n+1,
\]
and divides the chosen value by $2$ until it is odd. If the resulting odd value
is $1$, the player who made the move wins. Players alternate with perfect
information.

For the player to move, a position is \WIN\ if some move leads to a \LOSS\
position, and is \LOSS\ if every move leads to a \WIN\ position. Because the
game graph is infinite, these inductively generated classes need not a priori
exhaust all positions. Every remaining position is called \DRAW. This is the
standard finite-remoteness interpretation of Conway's Beans-Don't-Talk game
\cite{Guy1986,Guy1996}; the modern prize formulation is given by Alth\"ofer
\cite{AlthoferPage,Althofer2024}.

\begin{theorem}[Main theorem]\label{thm:main}
Every positive odd starting position has finite remoteness. Equivalently, the
game has no \DRAW\ positions, and optimal play reaches $1$ after finitely many
moves.
\end{theorem}

The proof is not a termination proof for arbitrary play. For example,
\[
  5\longrightarrow7\longrightarrow5
\]
is a legal two-ply cycle under suitable choices.

\section{Binary conjugation}

Write an odd original state as $n=2m+1$ and put
\[
  F(m)=\left\lceil\frac{3m}{2}\right\rceil.
\]
For a nonnegative integer $x$, let $R(x)$ be the integer obtained by deleting
the maximal alternating suffix of the binary expansion of $x$. For example,
$R(1001_2)=10_2$ and $R(10101_2)=0$.

\begin{proposition}[First normal form]\label{prop:firstnormal}
For $m>0$, the two children in $m$-coordinates are exactly
\[
  F(m),\qquad F(R(m)).
\]
\end{proposition}

\begin{proof}
We have
\[
  3(2m+1)-1=2(3m+1),\qquad
  3(2m+1)+1=2(3m+2).
\]
Exactly one of $3m+1$ and $3m+2$ is odd. In the other expression, each
additional factor of two removed corresponds to crossing one more alternating
bit of the binary suffix of $m$. The division stops at the first repeated
adjacent bit, which is precisely deletion of the maximal alternating suffix.
The remaining odd value, returned to $m$-coordinates, is $F(R(m))$.
\end{proof}

Every child in Proposition~\ref{prop:firstnormal} lies in the image of $F$.
Representing $F(q)$ by $q$ gives the conjugated game
\[
  A(q)=F(q)=\left\lceil\frac{3q}{2}\right\rceil,
  \qquad B(q)=R(F(q)),
\]
with terminal state $q=0$.

\begin{proposition}[Strict contraction]\label{prop:Bcontracts}
For every $q>0$,
\[
  B(q)<q.
\]
\end{proposition}

\begin{proof}
Deleting a nonempty binary suffix gives
\[
 B(q)\le \left\lfloor\frac{F(q)}2\right\rfloor
 \le \left\lfloor\frac{3q+1}{4}\right\rfloor<q
\]
for $q\ge2$, and $B(1)=0$ directly.
\end{proof}

The difficulty is therefore concentrated in the expanding branch $A$ and in
the unbounded length of the alternating suffix deleted by $B$.

\section{Constant-tail coordinates and source coordinates}

For positive odd $a$, $r\ge1$, and $e\in\{0,1\}$ define
\[
  Q_r^e(a)=a2^r-e.
\]
This is the canonical form of a positive binary word with a maximal constant
suffix of $r$ copies of the bit $e$.

\begin{lemma}[Long-tail recurrence]\label{lem:longtail}
For every $r\ge3$,
\[
 \boxed{
 A(Q_r^e(a))=Q_{r-1}^e(3a),\qquad
 B(Q_r^e(a))=Q_{r-2}^e(3a).}
\]
Moreover the two boundary states share their contracting child:
\[
 \boxed{B(Q_1^e(a))=B(Q_2^e(a)).}
\]
\end{lemma}

\begin{proof}
For $r\ge3$, the expanding child still ends in at least two equal bits, so
its maximal alternating suffix consists only of the final bit. The boundary
identity follows by comparing the binary words of $3a-e$ and $6a-e$; the
additional appended bit extends the same alternating suffix and therefore
deletes to the same prefix.
\end{proof}

Define the embedded three-free coefficient
\[
  J(s)=2A(s)+1.
\]
Every positive odd coefficient has a unique factorization
\[
  a=3^kJ(s),\qquad k\ge0,
\]
which defines its coefficient source $s$.

\begin{lemma}[Embedded source transition]\label{lem:source-transition}
For $s>0$, the two signed boundary transitions from the coefficient $J(s)$
select exactly the two ordinary children $A(s)$ and $B(s)$. More precisely,
there is a phase
\[
 \alpha(s)=1-\left(\left\lfloor\frac s2\right\rfloor\bmod2\right)
\]
such that
\[
 \odd(3J(s)+1-2\alpha(s))=J(A(s))
\]
with $2$-adic valuation one, while the opposite phase selects $J(B(s))$ with
valuation at least two.
\end{lemma}

The point of this embedding is that multiplication of a constant-tail
coefficient by $3$ preserves its source, while a signed exponent-one boundary
reveals one ordinary $A/B$ move on the source.

\section{Finite outcome certificates and token rank}

Every certified \WIN\ or \LOSS\ position has a finite canonical proof height:
\[
 h(0)=0,
\]
\[
 h(x)=1+\min\{h(y):y\text{ is a \LOSS\ child of }x\}
 \quad(x\text{ \WIN}),
\]
and
\[
 h(x)=1+\max\{h(y):y\text{ is a child of }x\}
 \quad(x\text{ \LOSS}).
\]

A marked normalization configuration carries only finite states whose outcome
and height comparison have actually been proved. These states are called proof
tokens.

\begin{lemma}[Multiset token rank]\label{lem:multiset}
For a finite multiset $M$ of proof heights define
\[
 \PhiM(M)=\mathop{\#}_{h\in M}\omega^h,
\]
the natural ordinal sum of the monomials $\omega^h$. Replacing one token of
height $H$ by any finite multiset of certified descendants of heights strictly
below $H$ strictly decreases $\PhiM$.
\end{lemma}

\begin{proof}
In Cantor normal form the coefficient of $\omega^H$ decreases by one, while
all inserted terms have lower exponent. No collection of lower monomials can
compensate for the loss at exponent $H$.
\end{proof}

This rank permits an exceptional local diamond to split one token into several
certified descendants without losing well-foundedness.

\section{Universal factor scan}

Fix a three-free odd base $a=J(s)$ and a phase $e$. For $k\ge0$ put
\[
 P_k=Q_2^e(3^ka),\qquad R_k=Q_1^e(3^ka),
\]
\[
 C_k=B(P_k)=B(R_k),\qquad T_k=A(R_k).
\]
Then
\[
 \operatorname{moves}(P_k)=\{R_{k+1},C_k\},\qquad
 \operatorname{moves}(R_k)=\{T_k,C_k\}.
\]

\begin{lemma}[Two-level DRAW horizon]\label{lem:twolevel}
If the adjacent frame $\{P_k,R_k\}$ contains a \DRAW, then
\[
 \boxed{\{C_k,T_k,C_{k+1},T_{k+1}\}\text{ contains a \DRAW}.}
\]
\end{lemma}

\begin{proof}
The common child $C_k$ cannot be \LOSS, otherwise both frame members would be
\WIN. If $C_k$ is \DRAW\ we are done. If $C_k$ is \WIN\ and $R_k$ is \DRAW,
then $T_k$ is \DRAW. The only remaining case has $P_k$ \DRAW\ and $R_k$
non-draw; then $R_{k+1}$ is \DRAW, so one of its two children
$T_{k+1},C_{k+1}$ is \DRAW.
\end{proof}

Factor at either exposed level
\[
 3^{i+1}J(s)+1-2e=2^vJ(t),\qquad v\ge1.
\]
Then
\[
 T_i=Q_v^{1-e}(J(t)).
\]
If $v\ge2$, the raw sibling is exactly
\[
 \boxed{C_i=Q_{v-1}^{1-e}(J(t)).}
\]
If $v=1$ and the positive raw sibling is \DRAW, its coefficient source is
strictly smaller than $t$. Thus every factorful exponent-one \DRAW\ reaches a
factor-free lift, a factor-free adjacent frame, or a genuine source descent
within at most two consecutive factor levels.

A second reverse argument removes the unbounded factor exponent itself. If a
factor frame does not pull back to a factor-free \DRAW\ over the same source,
then its only residual orientation is a marked gate
\[
  P_k\text{ \LOSS},\qquad R_k\text{ \DRAW}.
\]
The exact predecessor $Q_2^e(3^{k-1}J(s))$ is then either \DRAW, in which case
factors of three are removed at exponent two, or \WIN, in which case its
other child is an actual \LOSS\ proof token. The nonexceptional side exposes a
strictly lower \WIN\ descendant; the exceptional side is an exact adjacent
frame over that \LOSS\ token. Hence no factor level creates an unmarked reset.

\section{The typed fixed-fibre normalizer}

The local source, factor, and proof-token lemmas form a finite collection of
routing macros. Their full case proofs are recorded in Sections 91--135 of the
proof supplement at the repository revision cited in
Section~\ref{sec:repro}. We use only the following assembled statement.

\begin{proposition}[Typed normalizer]\label{prop:typed-normalizer}
Suppose a marked entry has already been supplied with
\begin{enumerate}[label=(\roman*)]
\item a retained arithmetic source anchor,
\item a finite multiset of certified \WIN/\LOSS\ proof tokens with proved
      provenance, and
\item one of the typed obligation/factor control states of the supplement.
\end{enumerate}
Then the equal-retained-data routing relation is well-founded.
\end{proposition}

The crucial scope restriction is that temporary routing cursors are not rank
anchors. In particular a canonical continuation
\[
 O(x,e)\longrightarrow O(A(x),1-e)
\]
usually increases the displayed ordinary source $x$, but this $A(x)$ is only
routing data. Sections 93--96 bound or rank the canonical $A$-streak and the
subsequent B-selecting transfers. Sections 94--95 rank factor forks, and
Sections 97--135 transport the finite tokens through every high-return and
long-tail recurrence. Every operation that changes a retained source or token
is supported by a strict comparison; pure routing preserves the retained
projection and is well-founded inside the fixed fibre.

\begin{lemma}[Well-founded fibre principle]\label{lem:fibre}
Let a transition relation have an ordinal-valued projection $\pi$. If
$C\to C'$ implies $\pi(C')\le\pi(C)$ and the relation restricted to every
fibre $\pi^{-1}(\{\alpha\})$ is well-founded, then the full relation is
well-founded.
\end{lemma}

\section{The external audit and the repaired entry}\label{sec:audit}

An external audit of an earlier version found two genuine defects.

First, the proof used
\[
 \text{smaller canonical coefficient}
 \Longrightarrow
 \text{smaller coefficient source},
\]
which is false when powers of three are present. The explicit counterexample
is
\[
 s=1,\qquad J(1)=5,\qquad a=3J(1)=15,\qquad e=1.
\]
The common boundary child is
\[
 R(3a-e)=R(44)=22.
\]
Its canonical coefficient is $11<15$, but $11=J(3)$, so the source has
increased from $1$ to $3$.

Second, the reverse-factor lemma was used at exponent one even though it is
valid only from exponent two. Hence the family
\[
 Q_1^e(3^kJ(s)),\qquad k>0,
\]
needed a separate entry proof. Section~\ref{sec:audit} gives that proof
without reinstating either invalid shortcut.

\subsection{One-shot initialization}

The key observation is that there are two logically different operations:
initializing the typed normalizer after the raw entry has been normalized,
and resetting a normalizer that is already running. Only the second requires
a preceding strict source/token decrease.

Refine the retained lexicographic rank by an entry bit
\[
 \eta\in\{1,0\},\qquad 1>0,
\]
placed after the outer source and outer token multiset:
\[
 \boxed{
 (\text{outer source},\ \PhiM(M),\ \eta,\
   \text{inner source/token rank},\ \text{finite control rank}).}
\]
Here $\eta=1$ means that the current outer fibre has not yet initialized its
typed inner normalizer, and $\eta=0$ means that Proposition
\ref{prop:typed-normalizer} is running.

The transition
\[
 \boxed{\eta:1\longrightarrow0}
\]
is strict independently of the numerical value assigned to the newly
initialized inner source. Thus the first returned source may be larger than
the original source. Within an unchanged outer source/token fibre, the bit can
never return from $0$ to $1$. A later raw reinitialization is allowed only
after an earlier outer source or outer proof-token component has strictly
decreased, exactly as permitted by Lemma~\ref{lem:fibre}. Therefore the entry
bit cannot hide a repeated source-growth cycle.

\subsection{The factor-free exponent-two bootstrap}

Let $x>0$, $a=J(x)$, and put
\[
 P=Q_1^e(a),\qquad U=Q_2^e(a),
\]
\[
 b=B(P)=B(U),\qquad F=A(U)=Q_1^e(3a),\qquad g=1-e.
\]
Suppose $U$ is \DRAW\ at the globally least outer DRAW source $x$. The common
child $b$ has smaller canonical coefficient than $J(x)$, hence its coefficient
source is below $x$; so $b$ is not \DRAW. A child of a \DRAW\ cannot be
\LOSS, therefore
\[
 \boxed{b\text{ is \WIN},\qquad F\text{ is \DRAW}.}
\]
The exact first-factor identity is
\[
 \boxed{9a+1-2e=2^vJ(b)}.
\]
It follows that
\[
 A(F)=Q_v^g(J(b)).
\]
We now split according to whether the input phase at $x$ is B-selecting or
A-selecting.

\subsubsection{B-selecting input}

Assume $e=1-\alpha(x)$. The first source-boundary numerator is divisible by
$4$. The consecutive numerator relation gives
\[
 9J(x)+1-2e
 =3(3J(x)+1-2e)-2(1-2e)
 \equiv2\pmod4,
\]
so
\[
 \boxed{v=1.}
\]
Writing $u=A(x)$ and using the parity that characterizes the B-selecting
phase gives the exact common endpoint
\[
 \boxed{b=3A(x)+1.}
\]
The two children of the \DRAW\ state $F$ are therefore
\[
 T=Q_1^g(J(b)),\qquad C=R(T).
\]
The signed-prefix identity identifies the raw state $C$ with the ordinary
child of $b$ selected by phase $g$:
\[
 \boxed{C\in\{A(b),B(b)\}.}
\]
At least one of $T,C$ is \DRAW.

If $T$ is \DRAW, it is the direct exponent-one typed obligation at source
$b$, and the one-shot bit initializes the normalizer.

If $C$ is \DRAW, the other ordinary child $\ell$ of the \WIN\ state $b$ is
\LOSS. For nonexceptional $b$, the ordinary side diamond exposes
$B(A(b))$ as a \WIN\ child of $\ell$ and as a child of the continuing
\DRAW, hence a strict proof-token/boundary descent. If $b$ is exceptional
and $C=A(b)$, the exceptional side formula gives an adjacent \DRAW\ frame over
the actual \LOSS\ token $\ell$.

The sole remaining exceptional orientation is $C=B(b)$ \DRAW. Since
$b=3A(x)+1$,
\[
 b\le\frac{9x+5}{2}.
\]
For exceptional $b$, the alternating suffix of $A(b)$ has length at least
three, whence
\[
 C=B(b)\le\frac{A(b)}8\le\frac{3b+1}{16}.
\]
If $\rho(C)$ is the coefficient source of $C$, then
\[
 \rho(C)<\frac C6
 \le\frac{3b+1}{96}
 \le\frac{27x+17}{192}<x.
\]
Thus this last row is a genuine strict outer-source transition.

\subsubsection{A-selecting input}

Assume $e=\alpha(x)$. Lemma~\ref{lem:source-transition} gives
\[
 3J(x)+1-2e=2J(A(x)),
\]
and the consecutive numerator identity implies
\[
 \boxed{v\ge2.}
\]
Retain
\[
 9J(x)+1-2e=2^vJ(b).
\]
If $v\ge4$, then
\[
 16J(b)\le9J(x)+1\le27x+19.
\]
Since $J(b)\ge3b+1$,
\[
 48b+16\le27x+19,
 \qquad
 \boxed{b\le\frac{9x+1}{16}<x.}
\]
Both children of $F$ have factor-free source $b$, so this is a strict
outer-source exit.

Only $v=2,3$ remain. If $v=2$, $F$ is a \DRAW\ parent whose children are
\[
 Q_2^g(J(b)),\qquad Q_1^g(J(b)),
\]
which is exactly the second form of the typed obligation $O(b,g)$.

If $v=3$, the children are the valuation-three frame
\[
 Q_3^g(J(b)),\qquad Q_2^g(J(b)).
\]
The constructor automatically supplies the congruence hypothesis needed by
the first-factor hidden-parent lemma. Reducing
\[
 9J(x)+1-2e=8J(b)
\]
modulo $3$ and using $g=1-e$ yields
\[
 \boxed{J(b)\equiv1+g\pmod3.}
\]
Thus that hidden-parent lemma is used only in the typed rows for which its
hypothesis is proved. The ensuing factor-frame normalization produces a
strict source edge, a strict boundary/token edge, or the already typed
obligation entry.

We have proved the following.

\begin{proposition}[Factor-free base-entry attachment]\label{prop:baseattach}
Every factor-free exponent-two \DRAW\ at a globally minimal outer source has
an outcome-compatible finite continuation that either strictly lowers the
retained outer source, strictly lowers a certified finite proof token, or
consumes the one-shot entry bit and enters the typed normalizer.
\end{proposition}

\subsection{Arbitrary factorful exponent one}

Return to
\[
 Q_1^e(3^kJ(s))\text{ \DRAW},\qquad k>0.
\]
The predecessor reduction described after Lemma~\ref{lem:twolevel} has two
outputs. Either it returns to a factor-free \DRAW\ over the same source, in
which case the factor-free analysis terminates at exponent one or at the
exponent-two bootstrap of Proposition~\ref{prop:baseattach}; or it exposes a
factor-free \DRAW\ lift/frame together with an actual finite \WIN/\LOSS\
descendant of the predecessor token. In the latter case Lemma
\ref{lem:multiset} supplies a strict token edge, after which Lemma
\ref{lem:fibre} permits the inner arithmetic task to reset.

Consequently:
\begin{proposition}[Repaired arbitrary exponent-one attachment]\label{prop:attachment}
Every arbitrary factorful exponent-one \DRAW\ has a finite outcome-compatible
normalization that either strictly lowers the retained source, strictly lowers
an already marked finite proof token, or consumes the one-shot entry bit and
enters the typed normalizer. No comparison $t<s$ is required for the temporary
returned inner source.
\end{proposition}

\section{Proof of the main theorem}

\begin{proof}[Proof of Theorem~\ref{thm:main}]
Assume that a \DRAW\ exists and choose the least coefficient source $s$ of a
\DRAW. The source-zero case is handled by the explicit zero-source
constant-tail normalization in the supplement. Assume $s>0$.

The constant-tail/source analysis and the factor scan reduce the continuing
actual \DRAW\ to a marked macro-configuration. By Proposition
\ref{prop:attachment}, every macrostep has one of three effects:
\begin{enumerate}[label=(\arabic*)]
\item the retained outer source strictly decreases;
\item at fixed outer source, the finite proof-token multiset strictly
      decreases under $\PhiM$;
\item with both outer components fixed, the one-shot bit changes from
      $\eta=1$ to $\eta=0$, after which the route lies in the typed fixed-fibre
      normalizer of Proposition~\ref{prop:typed-normalizer}.
\end{enumerate}
The first effect contradicts minimality when it occurs at the initial outer
source and, more generally, is a strict first lexicographic component. The
second is well-founded by Lemma~\ref{lem:multiset}. The third is strict at the
entry bit and cannot reset back to $1$ while the preceding components are
unchanged. Once $\eta=0$, Proposition~\ref{prop:typed-normalizer} gives
well-foundedness of the equal-retained-data fibre. Lemma~\ref{lem:fibre}
therefore makes the complete marked macro relation well-founded.

On the other hand, every \DRAW\ position has no \LOSS\ child and has at least
one \DRAW\ child. Hence from any marked \DRAW\ configuration one may choose an
outcome-compatible continuing \DRAW\ and repeat the finite normalization
macrostep forever. This would produce an infinite path in the well-founded
marked relation, a contradiction.

Therefore the conjugated game has no \DRAW\ positions. Conjugacy with the
original odd-integer game proves that every positive odd starting state has
finite remoteness and optimal play reaches $1$.
\end{proof}

\section{Verification boundary and reproducibility}\label{sec:repro}

The mathematical claim above is a human proof. The computational artifacts
have deliberately narrower roles.

\begin{itemize}
\item The repository's standard-library Python tests check exact arithmetic
      identities, finite outcome certificates, and finite regression ranges.
\item The symbolic global-routing JSON checks the declared typed control/rank
      assembly. It does not encode the new one-shot $\eta$ bridge and therefore
      remains a conditional machine check rather than an end-to-end proof.
\item The Lean project kernel-checks general metatheory and selected definitions,
      but does not kernel-check the complete concrete no-\DRAW\ theorem.
\end{itemize}

The corrected proof package used for this manuscript is the branch
\path{repair/althoefer-audit-gap} of
\url{https://github.com/Grisha-Pochuev/3n-plus-minus-1-game},
with the closure addendum \path{docs/althoefer-audit-closure.md}. At the time
this manuscript was prepared, the corresponding draft pull request was PR~\#1.
The earlier archived preprint is
\url{https://doi.org/10.5281/zenodo.21844684}.
That archived version predates the external audit and should not be used in
place of this corrected manuscript.

The companion file \texttt{verify\_closure.py} supplied with this manuscript
is a standalone arithmetic regression for the new repair identities. It has
no third-party dependencies and explicitly labels itself as supporting
verification, not as a proof.

\section{Audit checklist for the repaired bridge}

An independent reviewer should try to falsify the following points first:
\begin{enumerate}
\item no step identifies a smaller canonical coefficient with a smaller
      coefficient source;
\item the reverse-factor lemma is never used at exponent one;
\item the B-selecting exponent-two base row really has factor valuation one;
\item the exceptional raw B-child estimate gives a source strictly below the
      old source;
\item the A-selecting row with $v\ge4$ really has $b<x$;
\item the $v=3$ constructor really satisfies the additional hidden-parent
      congruence before that lemma is invoked;
\item the entry bit is placed after the outer source/token components and
      never changes $0\to1$ inside an unchanged outer fibre;
\item every later reset is preceded by a genuine strict retained source or
      proof-token decrease;
\item each macro follows an actual outcome-compatible continuing \DRAW, while
      finite proof tokens are kept separate from temporary arithmetic cursors.
\end{enumerate}

\section*{Acknowledgment of the audit}
The repaired version was prompted by questions from Ingo Alth\"ofer concerning
the former global entry argument. Those questions led to the explicit
counterexample to the old source-descent inference and to the separate proof of
the exponent-one entry used here.

\begin{thebibliography}{9}

\bibitem{Guy1986}
R. K. Guy,
\emph{John Isbell's Game of Beanstalk and John Conway's Game of Beans-Don't-Talk},
Mathematics Magazine 59(5) (1986), 259--269.
DOI: \url{https://doi.org/10.1080/0025570X.1986.11977258}.

\bibitem{Guy1996}
R. K. Guy,
\emph{Unsolved Problems in Combinatorial Games}, Problem 42,
in \emph{Games of No Chance}, MSRI Publications 29, Cambridge University
Press, 1996.

\bibitem{Althofer2024}
I. Alth\"ofer, M. Hartisch, T. Zipproth,
\emph{Analysis of a Collatz Game and Other Variants of the $3n+1$ Problem},
Advances in Computer Games (2024), 123--132.
DOI: \url{https://doi.org/10.1007/978-3-031-54968-7_11}.

\bibitem{AlthoferPage}
I. Alth\"ofer,
\emph{Collatz Problems: The Two-Player $3n{+}{-}1$ Game},
\url{https://althofer.de/collatz-prizes.html}, accessed August 2026.

\bibitem{PochuevRepo}
G. Pochuev,
\emph{Optimal $3n\mathbin{\pm}1$ Game: Proof and Verification Repository},
\url{https://github.com/Grisha-Pochuev/3n-plus-minus-1-game}, 2026.

\bibitem{PochuevZenodo}
G. Pochuev,
\emph{Optimal Play in the Two-Player $3n\mathbin{\pm}1$ Game},
Zenodo record \url{https://doi.org/10.5281/zenodo.21844684}, 2026.
The record predates the external audit discussed in this corrected manuscript.

\end{thebibliography}

\end{document}
